% Prova fictícia com enumerate aninhado (dados de teste)
\documentclass{article}
\usepackage{enumitem,graphicx}
\begin{document}
\section*{Lista fictícia}
Instruções que não são questão.
\begin{enumerate}
\item Calcule $\int_0^1 x\,dx$.
  \begin{enumerate}[label=(\alph*)]
    \item $0$
    \item $\frac12$ \correta
    \item $1$
  \end{enumerate}

\item Prove que existem infinitos primos.

  \textbf{Resolução:} Suponha que haja finitos primos $p_1,\dots,p_n$ e considere $p_1\cdots p_n+1$.

\item Julgue em verdadeiro ou falso:
  \begin{itemize}
    \item (V) Todo número par é divisível por $2$.
    \item (F) $1$ é primo.
  \end{itemize}

\item Observe a figura \includegraphics[width=0.4\textwidth]{figura_teste} e responda.
  \begin{enumerate}
    \item Sub-item longo que descreve uma parte da questão discursiva, com bastante texto para não parecer alternativa \begin{itemize}\item interno\end{itemize}
    \item Outra parte.
  \end{enumerate}
  Explique o raciocínio.
\end{enumerate}
\end{document}
